\documentclass[letterpaper,journal]{IEEEtran}
\usepackage[UTF8]{ctex}
\usepackage{amsmath,amssymb,amsfonts}
\usepackage{booktabs}
\usepackage{array}
\usepackage{graphicx}
\usepackage{xcolor}
\usepackage{url}
\usepackage{cite}
\usepackage[hidelinks]{hyperref}

\hyphenation{op-tical net-works semi-conduc-tor PI-JWM Air-Fog-Sim}
\emergencystretch=5em
\hfuzz=40pt
\renewcommand{\arraystretch}{1.08}
\setlength{\tabcolsep}{3.6pt}
\setlength{\dblfloatsep}{8pt plus 2pt minus 2pt}
\setlength{\dbltextfloatsep}{10pt plus 2pt minus 2pt}
\newcolumntype{L}[1]{>{\raggedright\arraybackslash}p{#1}}
\newcommand{\code}[1]{\texttt{#1}}
\newcommand{\pjwm}{PI-JWM}
\newcommand{\tabfontsize}{\footnotesize}
\newcommand{\numfmt}[1]{\texttt{#1}}
\newcommand{\widecaption}[1]{\caption{#1}}
\newcommand{\gmp}{\mathrm{GMP}}
\newcommand{\hist}{\mathcal{H}}
\title{\pjwm{}：面向车载空地边缘计算的物理--信息联合世界模型}
\author{Yuyao~Shen}

\begin{document}
\maketitle

\begin{abstract}
车载空地边缘计算场景中，车辆移动性、无人机覆盖、无线链路质量、任务队列和通信/计算资源调度相互耦合，使得传统的单一链路预测、静态资源分配或短视任务卸载方法难以刻画系统未来状态。本文提出 \pjwm{}（Physical-Information Joint World Model），面向 UAV-assisted vehicular fog/edge computing 建立物理网络与信息网络联合状态表示，并学习动作条件下的多步状态转移。该模型以物理图描述位置、距离、速度和覆盖关系，以信息图描述链路活动、传输速率、任务流和资源动作，并通过 action-conditioned latent rollout 同时预测节点、链路、任务和资源状态。进一步地，本文将训练好的 world model 固定为 rollout evaluator，用于评估 autonomous 策略器生成的未来 offload/RB/CPU/return 动作质量。基于 action-aligned 数据的实验表明，当给定 true-future/oracle-style 动作时，\pjwm{} 的 active-rate RMSE 可达到 100.140209；而 autonomous policy bridge 的最好基线为 217.237962，说明系统主要瓶颈已经从 world model 本体转移到未来动作值幅值的生成质量。该结果支持将 \pjwm{} 作为物理--信息联合数字孪生与调度策略评估的基础模型。
\end{abstract}

\begin{IEEEkeywords}
Physical-information joint world model, vehicular fog computing, UAV-assisted edge computing, action-conditioned rollout, network digital twin, scheduling policy evaluation.
\end{IEEEkeywords}

\section{引言}

UAV-assisted vehicular fog/edge computing 同时包含物理移动过程和信息任务过程。车辆由道路交通流驱动，UAV 和 RSU/edge 节点提供动态覆盖与计算能力，任务在本地、空中、边缘和云侧之间卸载、传输、计算和回传。系统性能不仅取决于某一条无线链路的瞬时速率，也取决于未来几个时隙中车辆位置、链路活动、任务队列和资源分配动作的连续演化。因此，一个可用于调度评估的模型需要回答的问题不是“当前链路速率是多少”，而是“在给定一串未来调度动作后，整个物理--信息系统会如何变化”。

现有研究通常从三个角度处理该问题。第一类方法关注无线资源管理或任务卸载优化，直接在当前观测上求解 offload、RB 或 CPU 分配；第二类方法关注交通或无线链路预测，强调时空相关性建模；第三类方法借鉴数字孪生或 world model 思想，把复杂系统写成可滚动预测的状态转移模型。单独使用这些方法会遇到一个共同困难：物理网络和信息网络之间不是松散相关，而是通过调度动作强耦合。车辆运动改变覆盖和 CSI，CSI 决定可用 rate，rate 影响任务传输，任务队列又反过来影响 offload/RB/CPU 动作。

本文的核心思路是把该场景建模为动作条件下的物理--信息联合 world model。\pjwm{} 不把仿真日志展平成普通表格输入，而是显式保留物理图、信息图、动作序列和任务状态：物理图承载移动性与几何关系，信息图承载链路、任务和资源关系，动作序列作为未来状态转移条件。训练完成后，world model 既可以在 true-future/oracle-style 动作下给出预测上界，也可以作为 frozen rollout evaluator 检验 autonomous 策略器生成动作的质量。

本文贡献可以概括为四点：第一，提出面向空地车载边缘计算的物理--信息联合状态表示，将移动、通信、计算和任务状态放入统一预测框架；第二，构建 action-conditioned latent rollout，使模型能够在未来动作条件下多步预测节点、链路和任务状态；第三，引入 active-rate RMSE、activity F1 和 link RMSE 等指标，区分“链路是否活跃”和“活跃链路速率幅值”两个难度层次；第四，建立 frozen rollout evaluator 与 autonomous policy bridge，揭示 autonomous setting 下的性能差距主要来自策略器生成的 action value/magnitude，尤其是 RB 资源幅值，而不是 world model 的预测上界不足。

\section{相关工作}

\subsection{车载空地边缘计算与任务卸载}

车载边缘计算和移动边缘计算的核心问题是把终端侧任务、无线接入、边缘计算和云侧资源放入同一个时延--能耗--可靠性约束中联合优化。MEC 的通信视角综述指出，任务卸载不仅是计算放置问题，还受到无线传输速率、队列状态、边缘服务器负载和用户移动性的共同影响 \cite{mao2017mecsurvey}。Vehicular edge computing 进一步引入车辆高速移动、道路拓扑、RSU 覆盖和交通流时变性，使得同一调度动作在不同时间和空间位置上的后果并不相同 \cite{raza2019vecsurvey}。UAV-assisted edge/fog computing 则增加了三维覆盖、空地链路和 UAV 机动性，相关工作通常围绕 UAV 轨迹、任务卸载和通信/计算资源分配进行联合优化 \cite{li2021uavedgerl}。

这些工作为本文场景提供了系统背景，但大多直接求解当前或短期调度动作，即给定当前状态后输出 offloading、RB 或 CPU allocation。它们通常不会显式学习“给定一串未来调度动作后，物理移动、无线链路、任务队列和计算状态如何共同演化”。AirFogSim 提供了 UAV-integrated vehicular fog computing 的轻量级参考仿真环境 \cite{airfogsim2025tmc}，但本文不把 AirFogSim 作为主框架，而是把它视为车辆、UAV、任务、链路和动作日志的数据来源。PI-JWM 的目标是从这些日志中抽象出可滚动预测的物理--信息联合状态转移模型。

\subsection{无线网络数字孪生与 world model}

无线网络数字孪生强调用可更新、可校准、可预测的模型映射真实网络状态，并支撑网络规划、监测、优化和闭环控制 \cite{khan2022dt}。面向 6G 的 wireless network digital twin 进一步强调生成式模型、在线数据驱动更新和网络智能体协同 \cite{tao2024wndt}。这类研究提供了“用模型复现网络并预测未来”的总体方向，但落到车载空地边缘计算时，还需要解决动作条件性：未来状态不仅由当前网络状态决定，也由未来 offload/RB/CPU/return 等调度动作决定。

World model 从动作条件状态转移角度刻画环境动力学。早期 World Models 使用隐变量模型学习环境表征，并在隐空间中做想象式 rollout \cite{ha2018worldmodels}；Dreamer 进一步把 latent dynamics 用于行为学习和长 horizon imagination \cite{hafner2019dreamer}。从控制视角看，这类模型接近部分可观测系统中的 belief-state transition \cite{kaelbling1998pomdp}，区别在于现代神经 world model 用数据驱动方式学习隐状态和转移函数。\pjwm{} 结合无线网络数字孪生和 world model：它不是只拟合静态网络快照，而是学习调度动作条件下的物理--信息联合状态转移。

\subsection{图神经网络与时空预测}

无线系统天然具有图结构：节点可以表示车辆、UAV、RSU、edge server 或 cloud，边可以表示覆盖、干扰、通信链路、任务流或资源分配关系。GNN for wireless communications 的综述系统说明了图结构对无线资源管理、功率控制、链路调度和网络优化的适配性 \cite{shen2023gnnwireless}。ENGNN 进一步强调边更新机制对 radio resource management 的重要性，因为无线网络中很多关键变量本身就在边上，例如链路质量、干扰强度、RB 分配和速率 \cite{engnn2024}。

时空图神经网络在交通预测和城市计算中已经形成成熟路线：STGCN 通过时间卷积和图卷积捕捉交通流的时空相关 \cite{yu2018stgcn}，后续 STGNN 综述总结了动态图结构、长期依赖和多源城市数据融合中的核心问题 \cite{jin2024stgnn}\cite{jin2024gnnts}。PI-JWM 借鉴这些思想，但不直接把交通预测模型搬到本问题中。本文采用双图视角：物理图表示位置、距离、速度和覆盖，信息图表示链路活动、任务状态和资源动作；二者通过 action-conditioned rollout 耦合。

\subsection{决策接口与动作生成}

当 world model 固定后，调度器需要生成 future action condition。传统 model predictive control 使用模型预测候选动作后果，再根据目标函数选择动作 \cite{camacho2013mpc}；offline RL 和 behavior support 方法强调，脱离数据支持域的动作可能导致模型评估不可靠，因此需要约束候选动作分布 \cite{fujimoto2019bcq}\cite{kostrikov2022iql}。Selective prediction 和 conformal risk control 则提供了“模型不确定或风险过高时拒绝/保守处理”的工具 \cite{geifman2019selectivenet}\cite{angelopoulos2022crc}。Decision-focused learning 关注预测模型和下游决策目标之间的一致性 \cite{mandi2024dfl}。

这些工作共同说明，PI-JWM 不应只追求单步预测误差最低，还需要让预测结果可用于调度动作评估。本文把动作生成拆成 activity/activation 和 value/magnitude 两层：前者决定哪些 edge-step 被激活，后者决定 RB/CPU 等动作幅值。该分解使 frozen world model 不只是预测器，也可以作为调度动作质量的诊断接口。

\section{系统模型与问题定义}

本节定义 PI-JWM 的建模对象、符号和预测任务。系统以离散时间步 $t=1,2,\ldots$ 推进，每个时间步对应一次日志采样和调度更新。实体集合 $V$ 包含车辆、UAV、RSU/edge、cloud 以及任务相关节点；边集合根据语义分为物理边和信息边。表 \ref{tab:notation} 总结主要符号。

\begin{table*}[!t]
\widecaption{PI-JWM 系统模型中的主要符号与含义。}
\label{tab:notation}
\centering
\tabfontsize
\begin{tabular}{L{0.14\textwidth}L{0.78\textwidth}}
\toprule
符号 & 含义 \\
\midrule
$t,k,H,K$ & $t$ 为当前离散时间步，$k$ 为未来预测步索引，$H$ 为历史窗口长度，$K$ 为 rollout 预测步长。 \\
$V$ & 系统实体集合，包括车辆、UAV、RSU/edge、cloud 和任务相关节点。 \\
$G_t^{phy}$ & 时刻 $t$ 的物理图，描述位置、距离、速度、高度差、邻近关系和覆盖关系。 \\
$G_t^{info}$ & 时刻 $t$ 的信息图，描述链路 activity/rate、任务流、队列状态和资源动作。 \\
$E_t^{phy},E_t^{info}$ & 物理边集合和信息边集合；同一实体对可以同时具有物理关系和信息关系，但语义不同。 \\
$X_t^{phy},X_t^{info}$ & 物理图和信息图的节点/边特征，例如位置、速度、CSI、任务积压、历史 rate 等。 \\
$B_t^{phy},B_t^{info}$ & 图结构或关系掩码，用于表示候选连接、可覆盖关系、可通信关系和有效 edge-step。 \\
$a_t$ & 调度动作，包含 offload、RB、CPU 和 return 等动作计数或资源幅值。 \\
$q_t$ & 任务状态，包含任务生成、等待、传输、计算、完成和回传相关状态。 \\
$z_t$ & PI-JWM 的隐状态，承载从历史观测和历史动作中提取的系统动力学信息。 \\
$\hat{x},\hat{y}^{act},\hat{r},\hat{q}$ & 节点状态预测、链路活跃性预测、链路速率预测和任务状态预测。 \\
\bottomrule
\end{tabular}
\end{table*}

系统在时刻 $t$ 的联合状态写为
\[
s_t=\{G_t^{phy},G_t^{info},a_t,q_t\},
\]
其中 $s_t$ 不是单一向量，而是由图结构、图特征、动作和任务状态组成的结构化状态。物理图和信息图分别为
\[
G_t^{phy}=(V,E_t^{phy},X_t^{phy},B_t^{phy}),
\]
\[
G_t^{info}=(V,E_t^{info},X_t^{info},B_t^{info}).
\]
$V$ 是实体集合；$E_t^{phy}$ 描述“物理上谁接近谁、谁可能覆盖谁”，典型特征包括距离、相对位置、高度差、速度差和覆盖掩码；$E_t^{info}$ 描述“信息上谁和谁发生任务或资源交互”，典型特征包括 link activity、link rate、offload count、RB allocation、CPU allocation、return count 和任务流。$X_t^{phy}$、$X_t^{info}$ 分别存放对应图上的节点/边特征，$B_t^{phy}$、$B_t^{info}$ 表示候选边或有效关系掩码。无线系统天然具有图结构，GNN 在无线资源管理中的优势已经由无线 GNN 综述和 edge-update GNN 工作系统说明 \cite{shen2023gnnwireless}\cite{engnn2024}；因此，\pjwm{} 不把仿真日志直接展平成 MLP 输入，而是保留节点、边和关系类型。

双图拆分的目的在于避免把两类机制混在同一个黑箱特征中。物理图负责刻画移动性、几何关系和覆盖条件；信息图负责刻画链路、任务和资源动作。两张图通过同一实体集合 $V$ 对齐，并通过未来动作 $a_t$ 耦合：例如车辆移动改变 CSI，CSI 影响 rate，rate 决定任务传输进度，而任务积压又影响下一步 offload/RB/CPU 动作。

在历史窗口 $H$ 和预测步长 $K$ 下，世界模型学习
\[
\mathcal{F}_\theta:\left(s_{t-H+1:t},a_{t-H+1:t+K-1}\right)\mapsto \hat{s}_{t+1:t+K}.
\]
该公式表示：模型输入过去 $H$ 步观测状态 $s_{t-H+1:t}$，以及从历史末端到未来 $K$ 步所需的动作条件 $a_{t-H+1:t+K-1}$，输出未来 $K$ 步状态预测 $\hat{s}_{t+1:t+K}$。这里必须把未来动作纳入输入，因为 PI-JWM 预测的不是自然演化的无动作环境，而是受调度器控制的网络系统。若未来动作错误，即使历史状态完全正确，未来 rate、任务队列和计算状态也会偏离。

该形式对应动作条件状态空间模型。隐状态建模借鉴部分可观测系统中的 belief-state 思想和现代 world model 的 latent transition 形式 \cite{kaelbling1998pomdp}\cite{ha2018worldmodels}\cite{hafner2019dreamer}：
\[
z_t=f_\theta(s_{t-H+1:t},a_{t-H+1:t-1}),
\]
\[
z_{t+k}=T_\theta(z_{t+k-1},a_{t+k-1}).
\]
第一式中，$f_\theta$ 是历史编码器，把历史状态和历史动作压缩成当前隐状态 $z_t$；第二式中，$T_\theta$ 是 transition module，用未来动作 $a_{t+k-1}$ 推进隐状态。这个递推结构使 PI-JWM 可以在同一历史状态下比较不同 future action sequence 的后果，从而成为策略器评估器。

解码器输出四类预测目标：
\[
\hat{x}_{t+k}=D_{node}(z_{t+k}),\quad
\hat{y}_{e,t+k}^{act}=D_{act}(z_{t+k},e),
\]
\[
\hat{r}_{e,t+k}=D_{rate}(z_{t+k},e),\quad
\hat{q}_{t+k}=D_{task}(z_{t+k}).
\]
$D_{node}$ 预测节点侧物理/系统状态，$D_{act}$ 预测边 $e$ 在未来步 $t+k$ 是否活跃，$D_{rate}$ 预测该边的通信速率幅值，$D_{task}$ 预测任务队列和任务阶段状态。这里把 activity 和 rate 分开是必要的：inactive edge 数量远大于 active edge，如果只优化全边 rate RMSE，模型可以通过把大量 inactive edge 预测为低速率获得较好平均误差，却无法解释真正发生通信/卸载的活跃链路。
\section{数据来源与动作变量}

本课题的数据来自参考仿真器和后处理日志，但建模对象不是仿真器本身，而是由日志抽象出来的物理--信息联合状态转移。参考仿真侧主要提供三类实体：第一类是由 SUMO 交通流产生的车辆轨迹和道路运动状态；第二类是 AirFogSim 场景中的 UAV、RSU/边缘节点和 cloud 等计算通信基础设施；第三类是任务、链路和资源分配记录。cloud 在日志中表示远端云计算节点或云侧计算资源池，通常作为卸载目的地之一，不参与道路物理运动。

AirFogSim 的仿真步内调度逻辑可以拆成若干 scheduler。\code{scheduleTraffic} 负责推进车辆交通状态，车辆到达、行驶、离开主要由 SUMO 路网、route/trip 配置、随机种子和交通需求共同决定；\code{scheduleMission} 负责任务/mission 层事件；\code{scheduleOffloading} 负责决定任务卸载到本地、UAV、RSU/边缘或 cloud；\code{scheduleCommunication} 负责无线链路、信道状态、RB 分配和传输速率；\code{scheduleComputing} 负责 CPU 资源分配和任务计算推进；\code{scheduleReturning} 负责计算结果回传路径和回传任务完成。UAV 运动不应简单写成“一定匀速直线”，更准确的说法是：仿真器根据 mobility/mission 配置更新 UAV 位置，具体是否近似匀速、是否按航点移动，取决于该次场景配置；PI-JWM 使用日志中的位置、速度、距离和覆盖关系作为观测状态，而不把某一种运动假设写死进模型。

任务生成由 \code{task\_profile} 等配置控制。\code{task\_profile} 表示节点任务生成配置，通常包括任务到达强度、数据量、计算需求、截止时间或优先级等字段；\code{lambda=0.5} 可理解为任务到达过程的强度参数，常用于泊松式或近似泊松式任务生成；\code{dag\_edge\_prob=0.0} 表示任务 DAG 内部依赖边生成概率为 0，即当前任务多按独立任务处理，没有额外子任务依赖。每个任务不是抽象标签，而是带有输入数据量、计算量、时延约束、源节点和卸载/计算/回传状态的工作负载。任务卸载具体指把这些计算任务从产生节点转移到本地以外的计算节点执行，并在需要时回传结果。

通信日志中的 \code{rate\_sum} 表示一个时间步或一个统计窗口内链路传输速率的聚合值，通常由活跃链路的瞬时 rate 求和得到；\code{csi\_mean} 表示 CSI/channel state information 的均值统计，反映平均信道质量。PI-JWM 预测的 rate 不是单纯为了拟合一个通信物理量，而是为了刻画未来动作条件下“任务能否被及时传输和计算”的关键中间状态：rate 影响传输时延、任务队列演化、卸载成功率和后续资源调度可行性。

动作变量在策略器中按 edge/sample/step 组织。本文主要使用六个动作维度：\code{offload\_count} 表示卸载任务数量，\code{rb\_task\_count} 表示被分配 RB 的任务数量，\code{rb\_total} 表示 RB 资源总量，\code{cpu\_task\_count} 表示被分配 CPU 的任务数量，\code{cpu\_total} 表示 CPU 资源总量，\code{return\_count} 表示回传任务数量。RB 是 radio resource block，只属于无线频谱/时频资源，不包含 CPU 和存储；CPU 是计算资源，日志中通常是仿真归一化或配置单位下的计算能力/计算量分配，不一定直接等同真实 Hz；存储不作为本文策略器的主要动作维度。资源需求由任务配置和调度器共同决定，PI-JWM 学习的是“给定历史状态与未来资源动作后，链路、任务和节点状态如何变化”。

因此，策略器的作用不是替代 world model，而是生成 world model 所需的 future action condition。可以把它理解为两个层次：第一层是 activity/activation，决定哪些 edge-step 上存在卸载、RB、CPU 或回传动作；第二层是 value/magnitude，决定这些被激活动作的数量或资源幅值。该分解使得调度器既能学习“哪些关系会发生资源交互”，也能学习“发生交互时资源幅值是多少”。

\section{PI-JWM 方法}

\subsection{物理--信息双图表示}

\pjwm{} 的核心表示由物理图和信息图组成。物理图编码车辆、UAV、RSU/edge 和 cloud 等实体之间的几何与覆盖关系，包括位置、距离、速度差、高度差和邻近关系；信息图编码任务、链路和资源动作，包括 link activity、rate、offload、RB、CPU、return 和任务队列状态。双图表示的目的不是增加模型复杂度，而是将两类因果链分开：物理图解释“谁能连接谁、连接条件如何变化”，信息图解释“连接上发生了什么任务和资源流动”。

\subsection{动作条件 latent rollout}

在历史窗口 $H$ 内，编码器从观测状态和历史动作中得到隐状态 $z_t$；在预测步 $k$ 上，transition module 使用未来动作 $a_{t+k-1}$ 推进隐状态：
\[
z_{t+k}=T_\theta(z_{t+k-1},a_{t+k-1}).
\]
这种写法把调度动作作为状态转移条件，而不是作为预测后的附加特征。其好处是可以比较不同 future action condition 下的系统演化，从而支持后续策略器评估和候选动作诊断。

\subsection{多任务解码器}

解码器分别输出节点状态、链路 activity、链路 rate 和任务状态。activity head 负责判断未来 edge-step 是否活跃；rate head 负责预测活跃链路上的速率幅值；task head 负责预测任务队列、完成或积压相关状态。由于 inactive edge 数量远多于 active edge，rate 建模不能只看全边平均误差，而必须单独报告 active-rate RMSE。该指标直接对应“真正发生通信/卸载的链路速率是否预测正确”。

\subsection{Frozen rollout evaluator 与策略器接口}

训练好的 \pjwm{} 可以被冻结为 rollout evaluator。给定同一个历史状态，不同策略器会生成不同的 future action tensor；frozen world model 将这些动作映射到未来状态预测，并用 active-rate RMSE、link RMSE、activity F1 和 task RMSE 等指标评价动作质量。这个接口把“学环境动力学”和“生成调度动作”分开：前者由 world model 完成，后者由 autonomous policy bridge 或后续调度器完成。

\section{训练目标与评价指标}

整体 link-rate RMSE 会被大量 inactive edge 稀释，因此本文同时报告 active-rate RMSE。真实活跃边集合定义为
\[
\Omega_{act}=\{(e,k)\mid y_{e,t+k}^{act}=1\}.
\]
active-rate RMSE 为
\[
\mathrm{RMSE}_{act-rate}
=\sqrt{\frac{1}{|\Omega_{act}|}\sum_{(e,k)\in\Omega_{act}}
(\hat r_{e,t+k}-r_{e,t+k})^2}.
\]
任务状态 RMSE 写为
\[
\mathrm{RMSE}_{task}
=\sqrt{\frac{1}{NKT}\sum_{n=1}^{N}\sum_{k=1}^{K}\sum_{j=1}^{T}
(\hat q_{n,k,j}-q_{n,k,j})^2}.
\]
activity 使用 precision、recall 和 F1：
\[
\mathrm{F1}=\frac{2\cdot \mathrm{Precision}\cdot \mathrm{Recall}}{\mathrm{Precision}+\mathrm{Recall}}.
\]

后续结果表中，$\uparrow$ 表示越大越好，$\downarrow$ 表示越小越好；粗体表示该表内该指标的最优值。

多任务训练目标为
\[
\mathcal{L}=\lambda_{node}\mathcal{L}_{node}
+\lambda_{act}\mathcal{L}_{act}
+\lambda_{rate}\mathcal{L}_{rate}
+\lambda_{task}\mathcal{L}_{task}.
\]
其中 activity 使用二分类交叉熵：
\[
\mathcal{L}_{act}=-y\log p-(1-y)\log(1-p).
\]
rate 使用 active-mixed loss：
\[
\mathcal{L}_{rate}=\mathcal{L}_{act\text{-}rate}
+\alpha\mathcal{L}_{inact\text{-}rate},
\]
\[
\mathcal{L}_{act\text{-}rate}
=\frac{1}{|\Omega_{act}|}\sum_{(e,k)\in\Omega_{act}}
(\hat r_{e,t+k}-r_{e,t+k})^2,
\]
\[
\mathcal{L}_{inact\text{-}rate}
=\frac{1}{|\Omega_{inact}|}\sum_{(e,k)\in\Omega_{inact}}
\hat r_{e,t+k}^{2}.
\]
多任务 loss balancing 直接对应多任务学习中损失尺度和梯度冲突问题 \cite{zhang2022mtl}\cite{kendall2018uncertainty}\cite{chen2018gradnorm}。本文实验采用的 balanced 配置为
\[
(\lambda_{node},\lambda_{act},\lambda_{rate},\lambda_{task})=(1.0,1.5,0.3,1.0).
\]

\subsection{能耗与逐步收益诊断指标}

状态预测指标不足以说明候选动作是否值得执行。为避免用单一综合 reward 掩盖任务收益与资源代价，本文额外记录 UAV 能耗分项和逐步任务收益。对 UAV $u$ 的第 $k$ 个仿真步，参考仿真器中的能耗写为
\begin{equation}
e_{u,k}=c_{fly}I^{fly}_{u,k}+c_{hover}I^{hover}_{u,k}
+c_{sense}n^{sense}_{u,k}+c_{rx}b^{rx}_{u,k}+c_{tx}b^{tx}_{u,k},
\label{eq:uav_step_energy}
\end{equation}
其中 $I^{fly}$ 和 $I^{hover}$ 为飞行/悬停指示量，$n^{sense}$ 为传感器使用数量，$b^{rx}$ 和 $b^{tx}$ 为当步收发数据量。系数 $c_{\cdot}$ 由 AirFogSim energy 配置给出；本实验报告 simulator energy unit，不将其解释为已经校准的真实焦耳。候选 $c$ 在 $H$ 步内的累计能耗为
\begin{equation}
E_H(c)=\sum_{k=0}^{H-1}\sum_{u\in\mathcal U_k}e_{u,k}.
\label{eq:candidate_energy}
\end{equation}
能耗记录需满足分项守恒和余额守恒：
\begin{equation}
\delta^{energy}_k=E^{before}_k-E^{after}_k-
\sum_{u\in\mathcal U_k}e_{u,k}=0.
\label{eq:energy_balance}
\end{equation}

逐步任务收益定义为
\begin{equation}
\begin{aligned}
u_k(c)={}&\Delta N^{done}_k(c)-\Delta N^{fail}_k(c)\\
&+\eta\log\!\left(1+\Delta B_k(c)\right),\quad \eta=0.01.
\end{aligned}
\label{eq:step_task_utility}
\end{equation}
其中 $\Delta B_k\geq 0$ 为 AirFogSim 全局累计通信量 \code{channel.data\_size} 的步间增量，不与 V2I/V2U/U2I 子项重复累加。候选累计收益及其相对同 seed、同决策时刻默认候选 $c_0$ 的配对效应为
\begin{equation}
U_H(c)=\sum_{k=0}^{H-1}u_k(c),\qquad
\Delta U_H(c)=U_H(c)-U_H(c_0).
\label{eq:paired_task_utility}
\end{equation}
跨 horizon 比较采用单位步效应 $\overline{\Delta U}_H(c)=\Delta U_H(c)/H$，避免仅因累计步数增加而放大数值。若做能耗权重敏感性分析，则使用
\begin{equation}
\widetilde U_{H,\lambda}(c)=U_H(c)-\lambda
\frac{E_H(c)-E_H(c_0)}{\max\{|E_H(c_0)|,\epsilon\}},
\label{eq:energy_aware_sensitivity}
\end{equation}
但 $\lambda$ 不根据 validation 或 test 表现事后选择；原始任务和能耗分项仍是主报告量。

设 $\mathcal G$ 为 seed--decision-time 决策组集合，$\mathcal C_g$ 为组 $g$ 内候选，则候选区分度定义为
\begin{equation}
\rho_H=\frac{1}{|\mathcal G|}\sum_{g\in\mathcal G}
\mathbb I\!\left[\max_{c\in\mathcal C_g}U_H(c)-\min_{c\in\mathcal C_g}U_H(c)>10^{-8}\right].
\label{eq:candidate_distinguishability}
\end{equation}
$\rho_H$ 衡量有多少决策组能为 selector 提供非平凡监督，不等同于策略性能。

\section{实验结果与分析}

实验分析围绕三个问题展开：第一，\pjwm{} 在给定可靠 future action condition 时是否能学习物理--信息联合状态转移；第二，当未来动作由 autonomous 策略器生成时，性能差距主要来自 world model 还是动作生成器；第三，候选动作在不同 horizon 下是否形成可辨识的任务收益和 UAV 能耗差异。所有表中 $\uparrow$ 表示越大越好，$\downarrow$ 表示越小越好；粗体表示该表内最优。

\begin{table*}[!t]
\centering
\tabfontsize
\begin{tabular}{lrrrl}
\toprule
动作条件 & active-rate RMSE$\downarrow$ & link RMSE$\downarrow$ & activity F1$\uparrow$ & 说明 \\
\midrule
true-future / oracle-style actions & \textbf{100.140209} & 5.377113 & 0.902439 & 可靠未来动作下的 world-model 上界 \\
zero future actions & 502.550213 & 12.440349 & 0.014320 & 去掉未来动作条件后明显退化 \\
zero history actions & 97.710620 & \textbf{5.342461} & \textbf{0.902439} & 历史动作不是主要瓶颈 \\
\bottomrule
\end{tabular}
\caption{Future-action condition 对 \pjwm{} rollout 的影响。}
\label{tab:future_action_condition}
\end{table*}

表 \ref{tab:future_action_condition} 表明，future action condition 是 \pjwm{} 能否预测链路活动和速率幅值的关键输入。给定 true-future/oracle-style actions 时，active-rate RMSE 达到 100.140209；将未来动作置零后，active-rate RMSE 上升到 502.550213，activity F1 降到 0.014320。这说明模型并非只依赖历史状态做惯性外推，而是在利用未来 offload/RB/CPU/return 等动作条件刻画状态转移。

\begin{table*}[!t]
\centering
\tabfontsize
\begin{tabular}{lrrrl}
\toprule
动作生成方式 & active-rate RMSE$\downarrow$ & link RMSE$\downarrow$ & activity F1$\uparrow$ & 诊断含义 \\
\midrule
true-future / oracle-style actions & \textbf{100.140209} & \textbf{5.377113} & \textbf{0.902439} & world model 参考上界 \\
raw threshold policy bridge & 280.321712 & 107.205085 & 0.026096 & 直接策略输出不足 \\
codebook policy bridge & 217.237962 & 81.257763 & 0.023506 & autonomous baseline \\
policy activity + true value & 122.205507 & 5.564188 & 0.884298 & activation 路径基本可用 \\
true activity + policy value & 314.673188 & 49.902519 & 0.232181 & value/magnitude 是主要瓶颈 \\
\bottomrule
\end{tabular}
\caption{Frozen \pjwm{} evaluator 下的 autonomous policy bridge 诊断。}
\label{tab:policy_bridge_diagnosis}
\end{table*}

表 \ref{tab:policy_bridge_diagnosis} 进一步把误差来源拆成 activation 和 value/magnitude 两部分。autonomous policy bridge 的最好 baseline 为 217.237962，仍明显弱于 oracle-style 100.140209。半 oracle 诊断显示，保留策略器 activation、替换真值 value 时可以达到 122.205507；相反，使用真值 activation、保留策略器 value 时退化到 314.673188。因此，当前主要限制不是 world model 无法预测状态转移，也不是 activation/ranking 完全失效，而是被激活动作的 RB/CPU 等资源幅值不够准确。

\begin{table*}[t]
\centering
\tabfontsize
\begin{tabular}{lrrrrl}
\toprule
策略器候选 & 候选数 & val RMSE$\downarrow$ & val$\rightarrow$test RMSE$\downarrow$ & test oracle RMSE$\downarrow$ & 诊断含义 \\
\midrule
pairwise quota 512 & 13 & 220.537045 & 235.081397 & 208.788876 & validation 选择不稳定 \\
val/test256 seed20260630 & 13 & 229.616248 & 216.897951 & 201.554773 & 小样本可低于 autonomous baseline \\
val/test256 seed20260631 & 13 & 223.966536 & 213.041887 & 203.444126 & 小样本最好可部署信号 \\
strict512 seed20260631 & 13 & 224.375085 & 236.977850 & 210.975399 & 512 confirmation 未通过 \\
expanded55 strict512 & 55 & 226.949165 & 235.423373 & 215.394402 & 候选空间有 headroom，但 selector 未稳定选准 \\
\bottomrule
\end{tabular}
\caption{v11 candidate actual-rollout template selector 诊断。所有候选均在 frozen \pjwm{} rollout evaluator 下评估，输出索引见 \code{code/artifacts/reports/ppt\_v11\_candidate\_20260629/experiment\_evidence\_index.json}。Oracle 行只表示候选集合内的事后选择上界，不能当作可部署策略器结果。}
\label{tab:v11_actual_rollout_selector}
\end{table*}

进一步的 v11 candidate 实验把策略器训练从直接行为克隆转向 actual-rollout candidate selection。表 \ref{tab:v11_actual_rollout_selector} 显示，val/test256 的小规模验证中，best-val 对应 test 可以达到 216.897951 和 213.041887，说明候选模板内确实存在接近或超过 autonomous baseline 217.237962 的信号。然而，在更严格的 512 样本确认中，13 候选和 expanded55 候选的 best-val 对应 test 分别退化到 236.977850 和 235.423373。expanded55 的 test sample-oracle 为 215.394402，说明候选空间本身已经有接近 baseline 的 headroom；但可部署 selector 不能稳定选择这些候选。因此，当前瓶颈不是 frozen world model 的预测上界。selector 泛化和候选监督覆盖仍需分开检查，不能直接把下一步收缩为继续更换排序模型。

\subsection{Rollout horizon 与动作阶段诊断}

本实验固定 seeds 0--4，每个 seed 选择三个决策点，每个决策点最多八个候选，覆盖 RB、offload target、mixed offload--RB、CPU scale 和 return route。三个 horizon 使用完全相同的 66 个候选。候选动作只在决策时刻注入一次，此后各候选均使用相同的 default scheduler。因此，$H=3,10,20$ 比较的是一次性动作扰动的延迟后果，而不是持续执行三种不同策略。

\begin{table*}[!t]
\centering
\tabfontsize
\begin{tabular}{rrrrrl}
\toprule
$H$ & seeds & step rows & candidates & non-trivial groups $\uparrow$ & quality audit \\
\midrule
3  & 5 & 198  & 66 & \textbf{7/15 (46.67\%)} & pass \\
10 & 5 & 660  & 66 & \textbf{7/15 (46.67\%)} & pass \\
20 & 5 & 1320 & 66 & \textbf{7/15 (46.67\%)} & pass \\
\bottomrule
\end{tabular}
\caption{Rollout-horizon 诊断。三行使用同一候选集合；quality audit 检查式 \eqref{eq:energy_balance}、式 \eqref{eq:paired_task_utility} 的可重建性、负能耗和动作注入有效性。结果来源为 \code{code/artifacts/reports/pi\_jwm\_energy\_reward\_horizon\_study\_20260713/}。}
\label{tab:horizon_stage_diagnostic}
\end{table*}

\begin{figure*}[!t]
\centering
\includegraphics[width=0.74\textwidth]{figs/pi_jwm_horizon_nontrivial_ratio.png}
\caption{式 \eqref{eq:candidate_distinguishability} 定义的候选区分度。三个 horizon 均为 7/15，延长观察窗口没有增加非平凡监督覆盖。}
\label{fig:horizon_nontrivial_ratio}
\end{figure*}

\begin{figure*}[!t]
\centering
\includegraphics[width=0.82\textwidth]{figs/pi_jwm_horizon_family_task_effect_per_step.png}
\caption{各动作族的单位步配对 task-utility 效应。该图描述单步扰动的暂态影响，不是跨 horizon 的策略性能排名。}
\label{fig:horizon_family_task_effect}
\end{figure*}

表 \ref{tab:horizon_stage_diagnostic} 给出三点结果。首先，$\rho_H$ 始终为 46.67\%，因此 horizon 过短不是候选监督稀疏的唯一原因。其次，CPU-scale 的 12 个候选和 return-route 的 4 个候选在三个窗口中均无可辨识 task-utility 或能耗差异；这只说明当前决策点、动作幅度和单步注入协议没有产生可测影响，不能外推为长期无效。最后，offload target 与 mixed offload--RB 的单位步平均效应在 $H=10$ 分别为 $-0.014543$ 和 $-0.015761$，到 $H=20$ 回到 $-0.000129$ 和 $-0.000737$。该非单调变化更符合中期传输/任务状态损失随后部分补偿，而不是 horizon 增长导致持续退化。

不同 horizon 下的累计能耗配对差异基本不变，这是协议所致：只有第一步动作不同，后续动作共享。由此，本实验支持“有效候选监督覆盖不足且 offload 效应具有暂态性”，但不支持“更长 horizon 已改善 selector”或“某个能耗权重已提高策略性能”。下一步应先构造持续多步候选序列，并增加真正处于 compute/return 敏感阶段的决策点；当 $\rho_H$ 和跨 seed 方向一致性提高后，再训练使用 PI-JWM forecast 与不确定性的 rollout-aware selector。

\subsection{阶段条件收益 LCB 与 selector 根因归因}

前述 CandidateSet ranker 在 validation 上全部 defer，单纯改变 defer 或 Pareto 阈值不能解释失败。为区分“机会检测失败”和“候选排序失败”，我们在冻结 world model 后对 schema-v6 候选标签进行两步审计。第一步用 selected-edge、pooled-interaction 和完整 edge-step interaction 特征检测候选是否存在非零 raw-SSE 改善机会；第二步按候选和 episode local step 估计跨训练 seed 的收益方向、均值、方差与支持数。episode phase 由 $\text{sample\_id}\bmod390$ 恢复，推理时不使用 seed 身份、真实 future state、counterfactual outcome 或 $\text{action\_applied}$ 标签。

最终的 phase-conditioned benefit LCB 不是新的 world model，而是一个可解释的 selector 基线。calibration seeds 44--49 只确定执行门槛，validation seeds 50--59 只做一次性评价。候选只有在正向比例、平均 raw-SSE benefit 和最小支持数同时达标，并且不被可观测 task--energy proxy Pareto 支配时才执行；否则回退到 ranked-allocation baseline。该规则可以把“为什么执行”分解为 phase、候选收益统计和安全门，而不是依赖不可解释的总分。

\begin{table*}[!t]
\centering
\tabfontsize
\begin{tabular}{lr}
\toprule
指标 & validation seeds 50--59 \\
\midrule
ranked baseline active-rate RMSE$\downarrow$ & 233.7162 \\
phase-conditioned benefit LCB RMSE$\downarrow$ & \textbf{207.5399} \\
相对 baseline 改善$\uparrow$ & 26.1763 \\
改善 seed 数/总 seed 数 & 10/10 \\
有效执行次数 & 65 \\
执行后正收益精度$\uparrow$ & 93.85\% \\
负选择率$\downarrow$ & 6.15\% \\
Pareto 违规数$\downarrow$ & 0 \\
activity F1 变化 & $-0.000086$ \\
link RMSE 相对变化 & $-0.5738\%$ \\
\bottomrule
\end{tabular}
\caption{阶段条件收益 LCB 的 validation 结果。该结果标记为 \texttt{diagnostic\_only}；matched test 18--19 与 external holdout 60--69 在配置冻结前均未访问。}
\label{tab:phase_benefit_lcb}
\end{table*}

该方法将 validation active-rate RMSE 从 233.7162 降至 207.5399，10 个 validation seed 全部改善，65 次有效执行中的正收益精度为 93.85\%，且没有 Pareto 违规。它达到 B 级区间 $[200,213.160874)$，但尚未达到预注册的 A 级门槛 $<200$，所以不能把它写成已经定型的部署策略。更重要的是，rank-only 归因的 validation RMSE 为 292.90--312.33，说明此前失败主要来自候选排序表示，而不仅是 defer 过于保守；schema-v6 交互审计的 opportunity PR-AUC 为 0.9495，但 sample rank Spearman 仅为 0.0832，支持“机会可辨识、收益排序困难”的判断。

该结果也限制了后续解释。phase context 和训练期 benefit 统计提供了稳定的 validation 信号，但不能替代 matched test 或 external holdout；candidate-specific experts、phase smoothing、estimator ensemble、residual routing、phase-restricted kNN 和无 phase token ranker 均未超过 exact phase table。因而本文把阶段条件收益 LCB 定位为可解释的诊断基线，并将持续多步候选序列、更多 compute/return 敏感决策点和独立外部验证列为后续工作，而不把本轮的 validation 改善表述为最终策略器突破。

\section{结论}

本文提出 \pjwm{}，面向车载空地边缘计算建立物理--信息联合 world model。该模型以物理图描述移动性、几何距离和覆盖关系，以信息图描述链路活动、任务状态和资源动作，并通过 action-conditioned latent rollout 预测未来节点、链路和任务状态。实验表明，\pjwm{} 在可靠 future action condition 下可以形成稳定的状态预测上界；当未来动作由 autonomous policy bridge 生成时，误差主要来自动作 value/magnitude，尤其是资源幅值生成，而不是物理--信息联合 world model 本身失效。能耗与 horizon 诊断进一步显示，当前候选库只有 7/15 个决策组提供非平凡监督，且一次性动作扰动在 $H=3,10,20$ 下的区分度不变。阶段条件收益 LCB 在未访问测试集的 validation 上将 active-rate RMSE 从 233.7162 降至 207.5399，并保持 93.85\% 的执行后正收益精度；该结果属于 B 级诊断信号而非最终部署结论。因而现阶段应继续提高候选序列和决策点的可辨识性，并在配置冻结后进行一次独立 matched/external 验证；本文不把 validation-only selector 或能耗权重敏感性结果写成已定型策略。该结果说明，\pjwm{} 可以作为面向空地车载边缘计算的可滚动预测模型和调度动作评估器，同时也为定位动作候选与策略器的误差来源提供受控接口。
\sloppy
\begin{thebibliography}{99}

\bibitem{mao2017mecsurvey}
Y. Mao, C. You, J. Zhang, K. Huang, and K. B. Letaief.
\newblock A Survey on Mobile Edge Computing: The Communication Perspective.
\newblock \textit{IEEE Communications Surveys \& Tutorials}, 19(4):2322--2358, 2017.
\newblock DOI: \href{https://doi.org/10.1109/COMST.2017.2745201}{10.1109/COMST.2017.2745201}.

\bibitem{raza2019vecsurvey}
S. Raza, S. Wang, M. Ahmed, and M. R. Anwar.
\newblock A Survey on Vehicular Edge Computing: Architecture, Applications, Technical Issues, and Future Directions.
\newblock \textit{Wireless Communications and Mobile Computing}, 2019:3159762, 2019.
\newblock DOI: \href{https://doi.org/10.1155/2019/3159762}{10.1155/2019/3159762}.

\bibitem{li2021uavedgerl}
S. Li, X. Hu, and Y. Du.
\newblock Deep Reinforcement Learning for Computation Offloading and Resource Allocation in Unmanned-Aerial-Vehicle Assisted Edge Computing.
\newblock \textit{Sensors}, 21(19):6499, 2021.
\newblock DOI: \href{https://doi.org/10.3390/s21196499}{10.3390/s21196499}.
\bibitem{shen2023gnnwireless}
Y. Shen, J. Zhang, S. H. Song, and K. B. Letaief.
\newblock Graph Neural Networks for Wireless Communications: From Theory to Practice.
\newblock \textit{IEEE Transactions on Wireless Communications}, 22(5):3554--3569, 2023.
\newblock DOI: \href{https://doi.org/10.1109/TWC.2022.3219840}{10.1109/TWC.2022.3219840}.

\bibitem{engnn2024}
Y. Wang, Y. Li, Q. Shi, and Y.-C. Wu.
\newblock ENGNN: A General Edge-Update Empowered GNN Architecture for Radio Resource Management in Wireless Networks.
\newblock \textit{IEEE Transactions on Wireless Communications}, 23(6):5330--5344, 2024.
\newblock DOI: \href{https://doi.org/10.1109/TWC.2023.3325735}{10.1109/TWC.2023.3325735}.

\bibitem{jin2024stgnn}
G. Jin, Y. Liang, Y. Fang, Z. Shao, J. Huang, J. Zhang, and Y. Zheng.
\newblock Spatio-Temporal Graph Neural Networks for Predictive Learning in Urban Computing: A Survey.
\newblock \textit{IEEE Transactions on Knowledge and Data Engineering}, 36(10):5388--5408, 2024.
\newblock DOI: \href{https://doi.org/10.1109/TKDE.2023.3333824}{10.1109/TKDE.2023.3333824}.

\bibitem{jin2024gnnts}
M. Jin, Y.-F. Li, et al.
\newblock A Survey on Graph Neural Networks for Time Series: Forecasting, Classification, Imputation, and Anomaly Detection.
\newblock \textit{IEEE Transactions on Pattern Analysis and Machine Intelligence}, 46(12):10466--10485, 2024.
\newblock IEEE Computer Society document: \url{https://www.computer.org/csdl/journal/tp/2024/12/10636792/1ZqVmrBY0Vy}.

\bibitem{khan2022dt}
L. U. Khan, Z. Han, W. Saad, E. Hossain, M. Guizani, and C. S. Hong.
\newblock Digital Twin of Wireless Systems: Overview, Taxonomy, Challenges, and Opportunities.
\newblock \textit{IEEE Communications Surveys \& Tutorials}, 24(4):2230--2254, 2022.
\newblock DOI: \href{https://doi.org/10.1109/COMST.2022.3198273}{10.1109/COMST.2022.3198273}.

\bibitem{tao2024wndt}
Z. Tao, W. Xu, Y. Huang, X. Wang, and X. You.
\newblock Wireless Network Digital Twin for 6G: Generative AI as a Key Enabler.
\newblock \textit{IEEE Wireless Communications}, 31(4):24--31, 2024.
\newblock DOI: \href{https://doi.org/10.1109/MWC.002.2300564}{10.1109/MWC.002.2300564}.

\bibitem{zhang2022mtl}
Y. Zhang and Q. Yang.
\newblock A Survey on Multi-Task Learning.
\newblock \textit{IEEE Transactions on Knowledge and Data Engineering}, 34(12):5586--5609, 2022.
\newblock DOI: \href{https://doi.org/10.1109/TKDE.2021.3070203}{10.1109/TKDE.2021.3070203}.

\bibitem{kendall2018uncertainty}
A. Kendall, Y. Gal, and R. Cipolla.
\newblock Multi-Task Learning Using Uncertainty to Weigh Losses for Scene Geometry and Semantics.
\newblock \textit{IEEE/CVF Conference on Computer Vision and Pattern Recognition}, 2018.
\newblock DOI: \href{https://doi.org/10.1109/CVPR.2018.00781}{10.1109/CVPR.2018.00781}.

\bibitem{chen2018gradnorm}
Z. Chen, V. Badrinarayanan, C.-Y. Lee, and A. Rabinovich.
\newblock GradNorm: Gradient Normalization for Adaptive Loss Balancing in Deep Multitask Networks.
\newblock \textit{International Conference on Machine Learning}, 2018.

\bibitem{geifman2019selectivenet}
Y. Geifman and R. El-Yaniv.
\newblock SelectiveNet: A Deep Neural Network with an Integrated Reject Option.
\newblock \textit{International Conference on Machine Learning}, 2019.
\newblock URL: \url{https://proceedings.mlr.press/v97/geifman19a.html}.

\bibitem{angelopoulos2022crc}
A. N. Angelopoulos, S. Bates, E. J. Cand\`es, M. I. Jordan, and L. Lei.
\newblock Conformal Risk Control.
\newblock arXiv:2208.02814, 2022.

\bibitem{fujimoto2019bcq}
S. Fujimoto, D. Meger, and D. Precup.
\newblock Off-Policy Deep Reinforcement Learning without Exploration.
\newblock \textit{International Conference on Machine Learning}, 2019.
\newblock URL: \url{https://arxiv.org/abs/1812.02900}.

\bibitem{kostrikov2022iql}
I. Kostrikov, A. Nair, and S. Levine.
\newblock Offline Reinforcement Learning with Implicit Q-Learning.
\newblock \textit{International Conference on Learning Representations}, 2022.
\newblock URL: \url{https://arxiv.org/abs/2110.06169}.

\bibitem{mandi2024dfl}
J. Mandi, J. Kotary, S. Berden, M. Mulamba, V. Bucarey, T. Guns, and F. Fioretto.
\newblock Decision-Focused Learning: Foundations, State of the Art, Benchmark and Future Opportunities.
\newblock \textit{Journal of Artificial Intelligence Research}, 80:1623--1701, 2024.
\newblock URL: \url{https://www.jair.org/index.php/jair/article/view/15320}.


\bibitem{kaelbling1998pomdp}
L. P. Kaelbling, M. L. Littman, and A. R. Cassandra.
\newblock Planning and Acting in Partially Observable Stochastic Domains.
\newblock \textit{Artificial Intelligence}, 101(1--2):99--134, 1998.
\newblock DOI: \href{https://doi.org/10.1016/S0004-3702(98)00023-X}{10.1016/S0004-3702(98)00023-X}.

\bibitem{ha2018worldmodels}
D. Ha and J. Schmidhuber.
\newblock World Models.
\newblock arXiv:1803.10122, 2018.

\bibitem{hafner2019dreamer}
D. Hafner, T. Lillicrap, J. Ba, and M. Norouzi.
\newblock Dream to Control: Learning Behaviors by Latent Imagination.
\newblock \textit{International Conference on Learning Representations}, 2020.

\bibitem{camacho2013mpc}
E. F. Camacho and C. Bordons.
\newblock Model Predictive Control.
\newblock Springer, 2013.

\bibitem{yu2018stgcn}
B. Yu, H. Yin, and Z. Zhu.
\newblock Spatio-Temporal Graph Convolutional Networks: A Deep Learning Framework for Traffic Forecasting.
\newblock \textit{IJCAI}, 2018.

\bibitem{airfogsim2025tmc}
Z. Wei, B. Li, R. Zhang, X. Cheng, and L. Yang.
\newblock AirFogSim: A Light-Weight and Modular Simulator for UAV-Integrated Vehicular Fog Computing.
\newblock \textit{IEEE Transactions on Mobile Computing}, 2025.
\newblock DOI: \href{https://doi.org/10.1109/TMC.2025.3641373}{10.1109/TMC.2025.3641373}.

\end{thebibliography}

\end{document}
